% =============================================================================
%  2-Desarrollo.tex — Desarrollo / Argumentación central
% =============================================================================

\section{Desarrollo y Avances de Investigación}
\label{sec:desarrollo-conceptual}

Durante los primeros días de abril de 2026, mi asesor \textit{David López Flores}
me recomendó adoptar el sistema de composición tipográfica \termino{LaTeX},
lo que benefició enormemente el avance al poder gestionar la evolución del documento formal
(tesis) mediante un sistema de control de versiones (\termino{Git} y \termino{GitHub}).
Todos estos avances pueden consultarse en el siguiente repositorio:
\url{https://github.com/galigaribaldi/TesisUrbanismo}, donde se encuentra una breve
descripción del tema. Asimismo, en la sección \textit{Releases}:
\url{https://github.com/galigaribaldi/TesisUrbanismo/releases}, se puede acceder
a la última versión completada y revisada del documento de investigación, actualmente
en su versión \texttt{v0.4.1}.


\subsection{Cambio en Enfoque y herramientas}
\label{ssec:cambio-enfoque}

Al principio de la investigación (semestre 2025-2) se planteó el tema del corredor anillar;
sin embargo, aún no se contaba con el conocimiento metodológico y sustantivo necesario
para desarrollarlo con rigor. En un inicio, la investigación se abordó desde el punto
de vista logístico. El cambio de paradigma surgió a principios de abril, cuando se tomó
la decisión de adoptar un enfoque matemático-urbano-computacional que combina el uso de
Grafos Matemáticos, indicadores urbanos y la implementación de sistemas computacionales
de autoría propia: \termino{Apímetro} (la primera API de datos abiertos del Sistema de
Transporte de la \zmvm) y \termino{VFTModel} (software libre para analizar las rutas
geográficas de cualquier sistema de transporte mediante GeoJSON y Grafos Matemáticos).

Todo esto con el objetivo de demostrar que la geometría del sistema de transporte
actual de la \zmvm requiere un cambio de paradigma hacia el modelo circular, siendo el
primer caso de estudio el del \termino{Anillo Periférico}.
En un inicio, debido al cambio de herramientas y enfoque, se contempló analizar únicamente
la sección sur de la Ciudad de México dada su gran extensión geográfica; sin embargo, esta
limitación desapareció con la implementación de las herramientas tecnológicas mencionadas.

\subsection{Git como registro de avances}
\label{ssec:git-avances}

Como parte del seminario de investigación y del cronograma de trabajo establecido,
se ha realizado una reestructuración y avance sustantivo en el documento de tesis
a lo largo del semestre 2026-2. A través del sistema de control de versiones
(\texttt{git}) del repositorio de la tesis, es posible rastrear con precisión
la evolución metodológica, estructural y de contenido desde el inicio del semestre
hasta la fecha de corte actual (mayo de 2026).

El Cuadro~\ref{tab:bitacora_git} presenta una bitácora cronológica de los hitos
de desarrollo registrados en el repositorio, agrupados por semana y etapa. Cada
entrada corresponde a una versión etiquetada (\textit{tag}) o a un conjunto de
\textit{commits} con impacto documentado en la estructura o el contenido de la tesis.

% ─────────────────────────────────────────────────────────────────────────────
%  TABLA 1: Bitácora cronológica por hitos de versión (git log)
% ─────────────────────────────────────────────────────────────────────────────
% Anchos ajustados para que suma + 2×\tabcolsep×4 cols ≤ \textwidth (155mm):
% 2.4+2.0+3.3+5.8 = 13.5cm → 135mm + 16.9mm padding = 151.9mm < 155mm
\begin{longtable}{>{\bfseries}p{2.4cm} p{2.0cm} p{3.3cm} p{5.8cm}}
    \caption{Bitácora cronológica de avances en el repositorio de tesis (Corte: 24 mayo 2026).}
    \label{tab:bitacora_git} \\
    \toprule
    \textbf{Versión / Hito} & \textbf{Fecha} & \textbf{Ámbito} & \textbf{Descripción del avance} \\
    \midrule
    \endfirsthead

    \multicolumn{4}{l}{\footnotesize\textit{(Continuación del Cuadro~\ref{tab:bitacora_git})}} \\
    \toprule
    \textbf{Versión / Hito} & \textbf{Fecha} & \textbf{Ámbito} & \textbf{Descripción del avance} \\
    \midrule
    \endhead

    \midrule
    \multicolumn{4}{r}{\footnotesize\textit{Continúa en la siguiente página}} \\
    \endfoot

    \bottomrule
    \endlastfoot

    % ── Semana 1 (10–11 abril) ─────────────────────────────────────────────
    \multicolumn{4}{l}{\small\color{ColorPrincipal}\textit{Semana 1 — Configuración inicial del repositorio (10--11 de abril)}} \\
    \addlinespace[2pt]

    Plantilla base
    & 10 abr 2026
    & Infraestructura
    & Creación de la plantilla \LaTeX{} UNAM con estructura de carpetas correcta
      (\texttt{1edd8b7}). \\
    \addlinespace[3pt]

    Introducción y referencias
    & 11 abr 2026
    & Cap.~1 y Bibliografía
    & Redacción de la introducción general y completado de primeras referencias
      bibliográficas (\texttt{3faa245}). \\
    \addlinespace[3pt]

    Marco teórico — borrador
    & 11 abr 2026
    & Cap.~2 y Cap.~3
    & Actualización de antecedentes, definiciones iniciales y paleta de comandos
      \LaTeX{} (\texttt{3c555a1, 7940ed7}). \\
    \addlinespace[8pt]

    % ── Semana 2 (14 abril) ────────────────────────────────────────────────
    \multicolumn{4}{l}{\small\color{ColorPrincipal}\textit{Semana 2 — Reestructuración y versión v0.3.0 (14 de abril)}} \\
    \addlinespace[2pt]

    \texttt{v0.3.0}
    & 14 abr 2026
    & Infraestructura
    & Integración de herramientas de proyecto: Makefile, GitHub Actions (CI/CD),
      Git LFS para datos geoespaciales, licencia CC BY-NC 4.0
      y plantillas de \textit{issues} (\texttt{8a1467b}). \\
    \addlinespace[3pt]

    \texttt{v0.3.1} — Recalibración ZMVM
    & 14 abr 2026
    & Cap.~1, Metodología
    & \textbf{Redefinición crítica}: ajuste sustantivo del alcance geográfico y
      demográfico de la investigación, acotándolo exclusivamente a la Zona
      Metropolitana del Valle de México (ZMVM). Eliminación de planteamientos
      duplicados y corrección de hipótesis y objetivos
      (\texttt{071ecc1, ff8f7df}). \\
    \addlinespace[8pt]

    % ── Semana 3 (22 abril) ────────────────────────────────────────────────
    \multicolumn{4}{l}{\small\color{ColorPrincipal}\textit{Semana 3 — Teoría de grafos (22 de abril)}} \\
    \addlinespace[2pt]

    \texttt{v0.3.2} — Grafos Matemáticos
    & 22 abr 2026
    & Cap.~3
    & Completado el subcapítulo de Grafos Matemáticos (secciones 3.2.0--3.2.6):
      fundamentos eulerianos, grafos computacionales con mapas de adyacencia,
      cálculo de caminos mínimos (\dijkstra, \floyd),
      integración con \gis{} y \gtfs{}, e implementaciones tecnológicas
      (\texttt{4de0d6a, 3284ca6}). \\
    \addlinespace[8pt]

    % ── Semana 4 (9 mayo) ─────────────────────────────────────────────────
    \multicolumn{4}{l}{\small\color{ColorPrincipal}\textit{Semana 4 — Inicio de capas de código (9 de mayo)}} \\
    \addlinespace[2pt]

    Apímetro — fase 1
    & 9 may 2026
    & Cap.~4
    & Primera versión documentada de las capas de código del Apímetro:
      extracción y procesamiento de datos \gtfs{} desde la \api{}
      de transporte público (\texttt{fe36c71}). \\
    \addlinespace[8pt]

    % ── Semana 5 (15–18 mayo) ─────────────────────────────────────────────
    \multicolumn{4}{l}{\small\color{ColorPrincipal}\textit{Semana 5 — Resultados e integración final (15--18 de mayo)}} \\
    \addlinespace[2pt]

    Apímetro terminado
    & 15 may 2026
    & Cap.~4
    & Cierre del capítulo de Apímetro con flujo completo de procesamiento
      documentado (\texttt{51a465e}). \\
    \addlinespace[3pt]

    Cap.~5 — Modelo VFT
    & 15 may 2026
    & Cap.~5
    & Estructuración del capítulo de resultados: estadísticas del grafo,
      tabla de indicadores, fases 1--3 del modelo VFT (impedancia, cobertura,
      fuerza capilar, factor de desviación, penalización por transferencia,
      \textit{Node Score}) (\texttt{89fa819, b252695}). \\
    \addlinespace[3pt]

    Resultados de cobertura y \textit{DF}
    & 17 may 2026
    & Cap.~5
    & Cálculo terminado de los indicadores de Cobertura de red y Factor de
      Desviación ($DI$); integración de figuras y tablas comparativas
      (\texttt{394a470}). \\
    \addlinespace[3pt]

    \texttt{v0.4.0} — Integración Cap. 4 y 5
    & 18 may 2026
    & Cap.~4, Cap.~5
    & Merge a \texttt{main} de los capítulos de Capas de Código y Resultados.
      Versión estable con todos los capítulos de análisis integrados
      (\texttt{a0ac6bc}). \\
    \addlinespace[3pt]

    \texttt{v0.4.1} — Migración LFS
    & 18 may 2026
    & Infraestructura
    & \textbf{Corrección técnica crítica}: las imágenes de los capítulos 3 y 4
      superaban el límite de almacenamiento de GitHub. Se migraron a
      Git LFS y se habilitó el soporte en el flujo de CI, restaurando
      la compilación automática del PDF (\texttt{5ede766}). \\

\end{longtable}

Información obtenida y optimizada por la herramienta nativa de \termino{Látex}: 
\termino{\textit{latexdiff}}.
% ─────────────────────────────────────────────────────────────────────────────
%  TABLA 2: Autoevaluación por capítulo
% ─────────────────────────────────────────────────────────────────────────────
\vspace{1em}

El Cuadro~\ref{tab:autoevaluacion} complementa la bitácora anterior con una
autoevaluación del estado actual de cada capítulo, identificando las debilidades
que deberán atenderse en la siguiente etapa del proceso de investigación.

\begin{table}[H]
    \centering
    \caption{Autoevaluación por capítulo: avances, redefiniciones y apartados débiles (Corte: mayo 2026).}
    \label{tab:autoevaluacion}
    \fontsize{9pt}{11.5pt}\selectfont
    \begin{tabularx}{\textwidth}{
        >{\hsize=0.55\hsize\raggedright\arraybackslash}X
        >{\hsize=1.15\hsize\raggedright\arraybackslash}X
        >{\hsize=1.15\hsize\raggedright\arraybackslash}X
        >{\hsize=1.15\hsize\raggedright\arraybackslash}X}
        \toprule
        \textbf{Capítulo} &
        \textbf{Avance concreto} &
        \textbf{Redefinición / Dificultad} &
        \textbf{Estado actual y debilidad} \\
        \midrule

        \textbf{Cap. 1}\newline Introducción y\newline Planteamiento &
        Consolidación del planteamiento, marco referencial e hipótesis. &
        \textit{Redefinición (v0.3.1):} Recalibración del alcance hacia la ZMVM;
        eliminación de planteamientos duplicados. &
        \textit{Estado:} Terminado al 100\%. \newline
        \textit{Pendiente:} Ajuste narrativo de cohesión tras redactar las conclusiones. \\
        \addlinespace\hline\addlinespace

        \textbf{Caps. 2 y 3}\newline Marco Teórico y\newline Matemático &
        Cap.~2 concluido. Completado subcapítulo de Grafos Matemáticos
        (3.2.0--3.2.6) con rigor analítico y notación formal. &
        \textit{Dificultad:} Articular el lenguaje sociológico del espacio urbano
        con el modelo matemático de grafos exigió múltiples iteraciones teóricas. &
        \textit{Estado:} Capítulos estructurados y cerrados. \newline
        \textit{Pendiente:} Afinar la transición narrativa entre el marco
        sociológico y el lenguaje formal matemático. \\
        \addlinespace\hline\addlinespace

        \textbf{Cap. 4}\newline Capas de Código\newline (Apímetro) &
        Documentación del flujo completo del Apímetro: extracción \gtfs{},
        construcción del grafo y cálculo de indicadores. &
        \textit{Dificultad técnica:} Peso de datos geoespaciales bloqueó la
        compilación; se requirió migración completa a Git LFS. &
        \textit{Estado:} Capítulo cerrado. \newline
        \textit{Pendiente:} Revisar consistencia de los bloques de código
        con la versión final del Apímetro. A su vez que pasar los bloques de código
        técnico a una sección llamada \termino{Anexos}.\\
        \addlinespace\hline\addlinespace

        \textbf{Cap. 5}\newline Resultados\newline Modelo VFT &
        Extracción de indicadores de Cobertura, Fuerza Capilar y Factor de
        Desviación ($DI$). Fases 1--2 del modelo VFT documentadas. &
        \textit{Pendiente de cálculo:} Fase 3 (\textit{Node Score}, $B(v)$, $\Delta E$)
        aún requiere completar el procesamiento computacional. &
        \textit{Estado:} Resultados parciales integrados. \newline
        \textit{Pendiente:} Interpretación cualitativa de los indicadores
        frente a las hipótesis espaciales originales. \\
        \bottomrule
    \end{tabularx}
\end{table}

\vspace{1em}
\textit{Nota metodológica:} Para respaldar esta autoevaluación y hacer tangibles
las redefiniciones estructurales, en la sección siguiente se presenta un fragmento
de evidencia generado con la herramienta \texttt{latexdiff} sobre el repositorio
de la tesis, permitiendo aislar visualmente (en azul las adiciones, en rojo las
supresiones) los cambios de texto ocurridos entre versiones.